\section{梯度流与小流时展开}
\label{sec:introGF}
\subsection{梯度流形式}
在下文中，我们在 $d$ 维欧氏时空中工作，$d=4-2\epsilon$。
梯度流通过引入一个额外的质量量纲为 $-2$ 的坐标 $t$（不要与闵氏时空的时间坐标混淆）——称为流时——把 $d$ 维 QCD 扩展为 $(d+1)$ 维场论 \cite{Luscher:2010iy, Luscher:2011bx}。在梯度流形式中，依赖流时的胶子场 $B_{\mu}^a(x, t)$ 与夸克场 $\chi^i_{\alpha}(x, t)$
由如下流方程确定 \cite{Luscher:2010iy, Luscher:2011bx,Luscher:2013cpa}：
\begin{eqnarray}
\label{eq:flowequations}
\partial_t B_{\mu} &=& \mathcal{D}_{\nu}G_{\nu\mu} +\kappa \mathcal{D}_{\mu}\partial_{\nu}B_{\nu},\\
\partial\chi &= & \mathcal{D}_{\mu}\mathcal{D}_{\mu}\chi -\kappa (\partial_{\mu}B_{\mu})\chi,\\
\partial\bar{\chi} &=& \bar{\chi}\overleftarrow{\mathcal{D}}_{\mu}\overleftarrow{\mathcal{D}}_{\mu} +\kappa\bar{\chi}\partial_{\mu}B_{\mu},
\end{eqnarray}
其中 $\kappa$ 是一个在物理观测量中消去的规范参数，且
\begin{eqnarray}
G_{\mu\nu} &= &\partial_{\mu}B_{\nu} -\partial_{\nu}B_{\mu} + \left[B_{\mu}, B_{\nu}\right],\\
\mathcal{D}_{\mu} &=&\partial_{\mu} + \left[B_{\mu}, .\right],
\end{eqnarray}
边界条件为
\begin{eqnarray}
\label{eq:boundary}
B_{\mu}(x; t=0) &=& gA_{\mu}(x),\\
\chi(x; t=0) &=& \psi(x),\\
\bar{\chi}(x; t=0) &=& \bar{\psi}(x).
\end{eqnarray}


在微扰梯度流中，人们通过给传播子添加依赖流时的指数因子、并引入额外的``流线''和``流顶点''来扩展常规的 QCD 费曼规则 \cite{Luscher:2011bx}；关于梯度流费曼规则及更高阶计算技术的更多细节见文献 \cite{Artz:2019bpr}。在以下微扰计算中，我们采用方便的选择 $\kappa=1$。

采用重正化的物理参数后，流场 $B_{\mu}$ 无需进一步重正化 \cite{Luscher:2010iy,Luscher:2011bx}，但流夸克场 $\chi$ 需要重正化 \cite{Luscher:2013cpa}：
\begin{eqnarray}
\chi_{0}(x; t) = Z_{\chi}^{\frac{1}{2}} \chi(x; t), \, \, \, \bar{\chi}_0(x; t) = Z_{\chi}^{\frac{1}{2}} \bar{\chi}(x; t), \, \, \,  Z_{\chi} = 1- \frac{\alpha_s}{4\pi}C_F \frac{3}{\epsilon_{\rm UV}} + O(\alpha_s^2),
\end{eqnarray}
其中 $C_F = (N_c^2-1)/(2N_c)$，$N_c=3$ 为色数。
为避免维数正则化中的 $Z_{\chi}$ 与格点正则化中的对应量之间匹配带来的复杂化，文献 \cite{Makino:2014taa} 引入了带圈（ringed）费米子场 $\mathring{\chi}(x,t)$、$\mathring{\bar{\chi}}(x,t)$：
\begin{eqnarray}
\mathring{\chi}(x,t) = \mathring{Z}_{\chi}^{\frac{1}{2}}\chi_{0}(x,t),\, \,  \, \mathring{\bar{\chi}}(x,t) = \mathring{Z}_{\chi}^{\frac{1}{2}}\bar{\chi}_{0}(x,t),
\end{eqnarray}
其中
\begin{eqnarray}
 \mathring{Z}_{\chi}=\frac{-2N_c}{(4\pi)^2t^2
\langle \bar{\chi}_{0}(x,t)\overleftrightarrow{\slashed{\mathcal{D}}}\chi_{0}(x,t)\rangle|_{m=0}},
\end{eqnarray}
这是单一夸克味的情形，其中 $\overleftrightarrow{\mathcal{D}}^{\mu} = {\mathcal{D}}^{\mu}  - \overleftarrow{\mathcal{D}}^{\mu}$。$\mathring{Z}_{\chi}$ 的次次领头阶（nnLO）结果见文献 \cite{Artz:2019bpr}。在次领头阶（NLO），重正化因子 $\mathring{Z}_{\chi}$ 由下式给出 \cite{Makino:2014taa}
\begin{eqnarray}
\label{eq:ringedfermionren}
\mathring{Z}_{\chi}(t, \mu) = (8\pi t)^{-\epsilon}\left\{1+  \frac{\alpha_s}{4\pi}C_F\left[\frac{3}{\epsilon_{\rm UV}} + 3\log\left(2\mu^2 te^{\gamma_E}\right)-\log(432)\right] \right\},
\end{eqnarray}
其中因子 $(8\pi t)^{-\epsilon}$ 的引入是为了平衡带圈费米子场的量纲：其量纲为 $3/2$，而流费米子场的量纲为 $(d-1)/2$。




%%%%%%%%%%%%%%%%%%%%%%%%%%%%%%%%%%%%%%%%%%
\subsection{小流时展开}
\label{subsec:smallflowexp}
在本工作中，我们考虑 eq.~(\ref{eq:currentoperator}) 与 eq.~(\ref{eq:currentoperatorspin}) 定义的、在流时 $t$ 处流动的局域流算符。只要不引起混淆，我们对流动与未流动的流算符使用相同的符号。就威尔逊线算符而言，我们在流时 $t$ 处流动外场和威尔逊线。具体地说，这意味着所有直接连接到``重''辅助场线上的胶子场都在流时 $t$ 流动，而``重''辅助场本身保持不流动。在小流时区域，我们对流动的流算符使用``短距离''OPE。我们的目标是导出重正化流动流算符与 $\overline{\rm MS}$ 重正化未流动流算符之间的匹配系数：
\begin{eqnarray}
\label{eq:smalltexpansion}
\mathcal{O}^{\rm R} (t)= c_{\mathcal{O}}(t, \mu)\mathcal{O}^{\overline{\rm MS}} (\mu) + O(t),
\end{eqnarray}
其中我们忽略了算符混合。


我们在夸克场的无质量极限下工作，并且在可行时把所有外标度取为零。
这一策略性选择使我们能够利用 HQET 匹配计算的成熟技术。具体来说，我们不需要计算未流动的算符，因为在维数正则化的微扰计算中它们在全阶都给出无标度积分，于是只需从重正化的流动算符中减去 IR 极点即可得到匹配系数 $c_{\mathcal{O}}(t, \mu)$。

类似于流算符的重正化，我们可以把局域流算符的匹配系数 $c_{\mathcal{O}}(t, \mu)$ 分解为场匹配系数与相应顶点匹配系数的乘积。对四个局域流算符 $\bar{\psi} h_v$、$gF^{\mu\nu}_{\|\perp}h_v$、$gF^{\mu\nu}_{\perp\perp}h_v$ 与 $g\bar{h}_vF^{\mu\nu}_{\perp\perp}h_v$，局域流算符匹配系数与相应的场、顶点匹配系数之间的关系为
\begin{subequations}
\label{eq:operatorrelations}
\begin{eqnarray}
c_{\psi h_v} (t,\mu) &=& \mathring{\zeta}_{\psi}(t,\mu)\zeta_{h_v}^F(t,\mu)\zeta_{\psi h_v}(t,\mu),\\
c_{\|\perp} (t,\mu)&=& \zeta_{A}(t,\mu)\zeta_{h_v}^A(t,\mu)\zeta_{\|\perp}(t,\mu),\\
c_{\perp\perp}(t,\mu) &=& \zeta_{A}(t,\mu)\zeta_{h_v}^A(t,\mu)\zeta_{\perp\perp}(t,\mu),\\
c_F(t,\mu) &=& \zeta_{A}(t,\mu)(\zeta^F_{h_v}(t,\mu))^2\zeta_F(t,\mu).
\end{eqnarray}
\end{subequations}
这里 $\mathring{\zeta}_{\psi}(t,\mu)$、$\zeta_{h_v}^F (\zeta_{h_v}^A)$ 和 $\zeta_{A}(t,\mu)$ 分别是带圈费米子场 $\mathring{\chi}$、基础（伴随）表示下的``重''辅助场 $h_v$ 以及流胶子场 $B$ 的匹配系数。此外，$\zeta_{\psi h_v}(t,\mu)$、$\zeta_{\|\perp}(t,\mu)$、$\zeta_{\perp\perp}(t,\mu)$ 与 $\zeta_F(t,\mu)$ 表示相应流动局域流算符顶点修正引起的匹配系数。我们在第 \ref{sec:oneloop} 节把这些匹配系数计算到单圈阶。

从我们在第 \ref{sec:EFT} 节关于威尔逊线算符可乘重正化的研究中得到的一个关键认识是：从梯度流方案过渡到 $\overline{\rm MS}$ 方案时，威尔逊线算符的匹配系数可以直接由相应局域流算符的匹配系数导出。这是因为在一维辅助场形式中，不同时空点上重正化局域流算符的组合除减除掉的线性发散外不引入额外的 UV 发散，同时也因为联系小流时极限下局域算符的 eq.~(\ref{eq:smalltexpansion})。因此，在小流时极限下不需要对这些不同时空点的组合局域流算符再做附加匹配。
